\section{Synthesis and Positioning of This Work}
\label{sec:gap_analysis}

Table \ref{tab:literature_synthesis} summarizes the identified research areas, existing limitations in the literature, and how this project directly addresses them through its modular architecture and experimental design.

\begin{table}[H]
\centering
\caption{Synthesis of Related Work and Project Positioning}
\label{tab:literature_synthesis}
\small
\begin{tabularx}{\textwidth}{>{\hsize=0.7\hsize}X >{\hsize=1.1\hsize}X >{\hsize=1.1\hsize}X >{\hsize=1.1\hsize}X}
\toprule
\textbf{Research Area} & \textbf{Existing Work / Focus} & \textbf{Limitation in Literature} & \textbf{How This Project Responds} \\
\midrule
Production rollup benchmarks & Focus on cost, proving, and deployment behavior (e.g., Chaliasos et al. \citet{chaliasos2024analyzing}). & Internals cannot be easily modified under identical workloads. & Provides configurable internal controls via a modular architecture (Sequencer, Executor, Prover, Submitter). \\
\addlinespace
Public rollup datasets & Focus on historical activity and fees. & Observational rather than experimental. & Contributes a replayable, synthetic benchmark harness for controlled configuration sweeps. \\
\addlinespace
Data Availability studies & Focus on DA problem and solution taxonomy (e.g., Saif et al. \citet{saif2024dataavailability}). & Often evaluated separately from live execution benchmarking. & Treats DA as an active experimental variable, showing empirically that Blob/EIP-4844 DA reduces publication cost. \\
\addlinespace
Proof systems & Focus on cryptographic proof construction. & Lack of system-level rollup trade-off analysis. & Employs proof-delay abstraction and isolated RISC0 profiling (demonstrating step-function behavior). \\
\addlinespace
Production documentation & Focus on operational architecture. & Optimized for deployment, reliability, and security, not for research isolation. & Employs a strict research-grade modular structure designed explicitly for telemetry extraction and bottleneck isolation. \\
\bottomrule
\end{tabularx}
\end{table}

In contrast to existing observational studies of production mainnets, our prototype prioritizes controlled experimentation, modularity, and observability. This approach provides an environment to conduct repeatable, isolated studies to observe exactly how individual design choices---such as batch size limits, priority scheduling heuristics (like TimeBoost), or DA mode selection---affect the throughput, latency, and cost behavior without the noise of production mainnet traffic. The findings in Chapter 5, including the measurable impact of larger batches on amortized cost versus latency, and the specific cycle limits of RISC0, are a direct result of this controlled methodology.
